
\subsubsection{Uncertainty Quantification for Language Models}

Uncertainty estimation in neural networks has a rich history \citep{gal2016dropout}, with recent adaptations for large language models. Semantic entropy \citep{kuhn2023semantic} extends classical entropy-based uncertainty to account for semantic equivalence---responses that differ in surface form but convey the same meaning are clustered together before computing entropy. The method uses bidirectional NLI to determine semantic equivalence, then computes entropy over the cluster distribution. Higher entropy indicates greater uncertainty and potential hallucination. Follow-up work demonstrated semantic entropy's effectiveness on detecting confabulations in long-form generation \citep{farquhar2024detecting}.

However, semantic entropy has been evaluated primarily on TruthfulQA and related factual QA tasks. The reliance on NLI-based clustering raises questions about transfer to benchmarks with different hallucination definitions.

\subsubsection{Self-Consistency Methods}

SelfCheckGPT \citep{manakul2023selfcheckgpt} measures surface-level agreement across multiple generations. The intuition is that factual knowledge, being reproducible, should yield consistent responses, while hallucinations vary. The method computes pairwise similarity (using BERTScore, n-gram overlap, or NLI-based comparison) and aggregates into a consistency score. SelfCheckGPT was evaluated on WikiBio hallucination detection.

The gap: SelfCheckGPT uses different benchmarks, different sample counts, and different base models than semantic entropy evaluations. Direct comparison of the two methods requires matching these variables.

\subsubsection{Confidence Calibration}

Contextual calibration \citep{zhao2021calibrate} adjusts model confidence using content-free inputs to estimate inherent biases. Temperature scaling and Platt scaling \citep{guo2017calibration} post-process confidence scores to improve calibration. While these methods primarily target classification tasks, they provide potential baselines for hallucination detection.

\subsubsection{Hallucination Benchmarks}

TruthfulQA \citep{lin2022truthfulqa} evaluates model truthfulness on questions designed to elicit common misconceptions. The benchmark labels responses as truthful or not based on ''Best Answer'' matching. HaluEval \citep{li2023halueval} provides hallucinated and correct response pairs across QA, summarization, and dialogue tasks, with explicit hallucination labels. These benchmarks encode different notions of ''hallucination.''
